Freddie Mercury ha passat a la història com una de les millors veus del rock. La seva màgica veu ha estat objecte de discussió i estudi, també per a la ciència. El biofísic austríac Christian Herbst va estudiar la veu del cantant de Queen i va determinar que Mercury era un baríton amb un registre vocal que anava del fa 2 (al voltant de 92,2~Hz) al sol 5 (al voltant de 784~Hz).

\figura[0.5]{fig1.png}
{\centering\small\textsc{Font:} \textit{https://queenphotos.wordpress.com}.\par}

\begin{apartat}{1,25}
Calculeu les longituds d'ona dels sons més greus i més aguts que Mercury podia emetre.
\end{apartat}

\begin{apartat}{1,25}
L'any 1985, Queen va actuar al festival Rock in Rio, en un concert que va aplegar unes 350\,000 persones. En un moment de molta emoció, els assistents van començar a cantar \textit{a cappella} la famosa cançó \textit{Love of my life}. Si cada assistent al concert cantava amb una potència de $10^{-7}$~W, quin nivell d'intensitat sonora (en decibels) es podia percebre a 1~km del concert? (A aquesta distància, podeu considerar que el concert és una font puntual de so.)
\end{apartat}

\dades{%
  $I_0 = 10^{-12}\ \mathrm{W\,m^{-2}}$.\\
  La velocitat del so en l'aire és de $340\ \mathrm{m\,s^{-1}}$.}
